% Generated draft: requires mathematical and source review.
\documentclass[11pt]{article}
\usepackage[margin=25mm]{geometry}
\usepackage{fontspec}
\setmainfont{Latin Modern Roman}
\usepackage{amsmath,amssymb,mathtools,unicode-math}
\setmathfont{Latin Modern Math}
\usepackage{enumitem,needspace,xurl,hyperref}
\setlength{\parindent}{0pt}
\setlength{\parskip}{5pt}
\setlength{\emergencystretch}{3em}
\newcommand{\N}{\mathbb{N}}
\newcommand{\Z}{\mathbb{Z}}
\newcommand{\Q}{\mathbb{Q}}
\newcommand{\R}{\mathbb{R}}
\newcommand{\C}{\mathbb{C}}
\newcommand{\F}{\mathbb{F}}
\newcommand{\A}{\mathbb{A}}
\newcommand{\PP}{\mathbb{P}}
\newcommand{\E}{\mathbb{E}}
\newcommand{\eps}{\varepsilon}
\newcommand{\dd}{\,\mathrm{d}}
\begin{document}
\section*{UnsolvedMath ID: 30004300}
\textbf{Source problem code:} catalogue entry
\subsection{Definitions and mathematical statement}
Let \((a_n)_{n\geq 1}\) be a sequence of real numbers satisfying
\[
1<a_1<a_2<\cdots,
\qquad r_n:=\frac{a_{n+1}}{a_n}.
\]
The source assumes that the successive ratios \((r_n)_{n\geq 1}\) are increasing, but does not specify whether this means \(r_{n+1}>r_n\) or \(r_{n+1}\geq r_n\) for every \(n\geq 1\). This distinction remains unresolved.

The source asks what can be said about the number of solutions of
\[
a_{i_1}+\cdots+a_{i_k}=a_{j_1}+\cdots+a_{j_k}.
\]
It does not specify a finite index range, the quantification or range of \(k\), whether repetitions of indices are allowed, whether permutations count separately, or whether solutions with identical multisets of indices on the two sides should be excluded.

To describe one possible finite counting formulation explicitly, for positive integers \(N\) and \(k\) define
\[
E_k(N;a):=
\#\left\{
(i_1,\ldots,i_k,j_1,\ldots,j_k)\in\{1,\ldots,N\}^{2k}:
\sum_{\ell=1}^{k}a_{i_\ell}=\sum_{\ell=1}^{k}a_{j_\ell}
\right\}.
\]
This definition counts ordered tuples, permits repeated indices, and includes solutions whose two index multisets coincide. It is an explicit auxiliary convention, not a convention established by the supplied source.

Under this convention, the source's question can be expressed as asking for estimates for \(E_k(N;a)\) under the stated sequence and ratio hypotheses. The source gives no particular bound or asymptotic assertion to prove and does not specify the desired dependence on \(N\), \(k\), or the sequence. Consequently, the intended counting problem requires clarification before a unique precise conjecture can be stated.
\subsection{Short English statement}
The problem asks how often two sums of \(k\) terms of a strictly increasing real sequence can agree when successive ratios increase. The meaning of increasing ratios, the index range, the counting conventions, and the desired estimate are unspecified. The displayed finite count provides one explicitly qualified formulation for review.
\subsection{Sources}
\begin{itemize}
\item Local UnsolvedMath registry, supplied record with permanent ID \(30004300\), title \emph{Equal Sum Counts for Increasing Ratio Sequences}, exact source problem code \texttt{catalogue entry}, in the Oberwolfach Reports Open Problems collection. No individual proposer is identified. The supplied proposed year is 2019.
\item Supplied report citation: \emph{Analytic Number Theory (2020)}, Oberwolfach Reports, DOI 10.4171/owr/2019/50, pp.\ 3141--3205. Source URL and supplied literature reference labeled ``original OWR report'': \url{https://doi.org/10.4171/owr/2019/50}. The supplied background describes Oberwolfach Reports as documenting mathematical workshops and their open-problem discussions. No problem-specific page or named contributor is supplied.
\item Literature assessment dated exactly \texttt{2026-08-22}: the supplied curation labels the problem open and states that known lacunary and geometric-ratio methods do not cover increasing, nonconstant ratios in the equal-\(k\)-sums problem, and that no general estimate for the size of such a set is established. This is supplied dated curation, not an independently verified assessment. The background expressly calls for independent verification of time-sensitive status and cited claims.
\item The supplied statement status is \texttt{corrected\_verified}. Its accompanying qualification says that only formula delimiters, notation, grammar, and typography were repaired, preserving mathematical content and scope. This supplied label does not constitute independent mathematical or source verification of this draft.
\end{itemize}
\end{document}
